% Auto-generated by scripts/build_judge_compare.py
\begin{table}[t]
\centering
\small
\setlength{\tabcolsep}{4pt}
\begin{tabular}{@{}lcccccc@{}}
\toprule
\textbf{Dimension} & \textbf{Ent} & \textbf{Cat} & \textbf{Sent} & \textbf{Span} & \textbf{Demo} & \textbf{Strict} \\
\midrule
Flavor & 40.0 & 40.0 & 80.0 & 93.3 & 36.7 & 33.3 \\
Texture & 96.7 & 93.3 & 90.0 & 100.0 & 93.3 & 83.3 \\
Meal Type & 66.7 & 70.0 & 66.7 & 96.7 & 63.3 & 46.7 \\
Occasion & 56.7 & 56.7 & 86.7 & 100.0 & 56.7 & 53.3 \\
Audience & 63.3 & 63.3 & 83.3 & 96.7 & 56.7 & 53.3 \\
Cuisine & 46.7 & 46.7 & 73.3 & 100.0 & 46.7 & 33.3 \\
\midrule
\textbf{Macro} & \textbf{61.7} & \textbf{61.7} & \textbf{80.0} & \textbf{97.8} & \textbf{58.9} & \textbf{50.6} \\
\bottomrule
\end{tabular}
\caption{Per-dimension breakdown of v2 judged precision ($n=30$ each). \emph{Demo} is the demo-relevant precision (entity, category, and span correct). Span groundedness is at or near 100\% on every dimension, supporting the anti-hallucination claim. Texture reaches near-publication quality (\textbf{93\% demo-relevant}, 83\% strict). The lowest scores live in Flavor and Cuisine, where the residual errors are structural: long-tail generic positive adjectives for Flavor, and modifier-vs-dish ambiguity (``Italian breadcrumbs'' $\not=$ Italian cuisine of the dish) for Cuisine. Both are extraction-prompt issues that the blacklist alone cannot address.}
\label{tab:llm-judge-perdim-v2}
\end{table}
